\section{Dataset Schema and Views}
\label{sec:schema-views}

The canonical candidate schema currently consists of decision metadata, candidate identity, split information, labels, and feature columns. The decision metadata includes \texttt{trace\_name}, \texttt{trace\_family}, \texttt{dataset\_source}, \texttt{capacity}, \texttt{horizon}, \texttt{decision\_id}, \texttt{decision\_t}, \texttt{decision\_chunk\_id}, and \texttt{split}. Candidate identity is represented through \texttt{candidate\_page\_id}. The labels are \texttt{y\_loss} and \texttt{y\_value}. Feature columns cover request context, candidate context, predictor/LRU signals, cache summary statistics, and recent request or hit rates.

Beyond candidate rows, the release defines two derived views. The \emph{decision view} summarizes each decision with candidate count, minimum loss, maximum value, the set of optimal candidate IDs, tie count, and regret summary statistics. The \emph{pairwise view} compares candidate pairs from the same decision and exposes pairwise preference labels. The current preserved open release does not ship the full pairwise view; it ships a capped pairwise sample of 1,000,000 rows instead.

\begin{figure}[t]
  \centering
  \fbox{\parbox{0.9\columnwidth}{\centering Placeholder for schema-and-view diagram:\\candidate schema, decision view, pairwise sample}}
  \caption{Planned schema and derived-view diagram.}
  \label{fig:schema-views}
\end{figure}

\begin{table}[t]
  \caption{Placeholder for compact schema summary.}
  \label{tab:schema-summary}
  \centering
  \begin{tabular}{ll}
    \toprule
    Column group & Notes \\
    \midrule
    Decision metadata & family, capacity, horizon, split, IDs \\
    Candidate identity & candidate page ID \\
    Labels & \texttt{y\_loss}, \texttt{y\_value} \\
    Feature groups & request, candidate, cache, disagreement, rate \\
    Derived decision fields & tie and regret summaries \\
    Derived pairwise fields & A/B labels and regret differences \\
    \bottomrule
  \end{tabular}
\end{table}
